\section{Discussion}\label{sec:discussion}

The strongest evidence for \builtnet{} is that the same input supports reproducible and interpretable rule changes. Field-level tests confirm that variants resolve as intended. Stable row indices then let the analyst compare a degree or closure statistic on aligned records. The tool's value is greatest when ambiguity is local. A choice about the Danish sampler can be represented by editing one field, so the consequences for the workplace layer can be isolated while the remaining code and inputs stay fixed. For the yearly Dutch and Swedish cases several release features differ together. A whole-card comparison answers a broader question about the composite specification.

The most obvious objection is that an approximate implementation cannot validate the actual register networks. The tests establish that the software executes its stated rules. The common-input comparison establishes the consequences of those operators for the given memberships. A comparison within the protected source environment would be needed to establish fidelity for release-specific dyads. Public aggregates can still reject some combinations of inputs and rules. A two-way draw produced a substantially closer Danish edge count, while its implied minimum degree conflicts with the source description. The Dutch secondary-school result similarly flags a mismatch between the open grouping labels and the target layer. In both cases a credible explanation remains conditional on how the original groups were formed.

The package supports a sequence of checks. First, compare core membership definitions and eligibility before projecting. Second, contrast a nominated-pair rule with a realised-degree cap using the same group assignments. Third, report which graph representation enters each statistic. Where group sizes are available without individual rows, an edge-free expectation gives a reference against which a result can be checked. Nearest-home closure and neighbourhood expansion need spatial information. The compact-population error for nearest-home clustering, the underprediction in the flat-heavy neighbourhood and the failed analytic correction for the yearly Dutch pair show how local assumptions propagate to comparison statistics.

The excess-closure counterexample also changes how to read a successful software cross-check. Agreement with igraph for ordinary triangles and with the comparison library for its weighted convention verifies those computations. It cannot decide whether the resulting statistic counts people, tie pairs or paths through layers. A release or paper should identify the numerator, denominator, edge union and treatment of undefined values. This matters for the README output as well. Its ordinary mean can become undefined because one active ego lacks a pair of alters. Reporting the count of defined coefficients makes the denominator inspectable.

There are two companion papers. One crosses populations and construction rules to study the comparability of the resulting statistics \citep{almamunbuilt2026}. The other develops edge-free formulas, investigates published aggregates and compares construction with observed contacts \citep{almamunpredicting2026}. Here, their results explain selected choices in the released package and the limits of public aggregate checks. The interface lets a researcher trace those choices from source statement to operator and from operator to statistic.

\section{Limitations}

\begin{itemize}
\item \textbf{Release coverage.} The cards implement a subset of each release's choices. Omitting an added relation can lower edge counts. Substituting a household for an address or retaining a broader school group can change degree in either direction. The direction for composite measures is unknown.
\item \textbf{Employment input.} The adapter has one employer and site per row. Omitting additional paid jobs generally removes opportunities arising from those jobs. If the comparison includes a new cap or uses an employer as a site, the net direction is unknown.
\item \textbf{Eligibility and missingness.} The Swedish age gate changes the denominator. Missing group data can remove edges from a layer, while permissive adapter defaults can create candidate pairs. The direction of a difference in mean clustering is unknown.
\item \textbf{Undirected construction.} The union guarantees at least as many ties as a mutual-nomination rule on the same nominations. It loses direction and can increase degree beyond the nominal outgoing cap. The change in local clustering has unknown direction.
\item \textbf{Spatial prediction.} A homogeneous household process lacks multihousehold addresses. The available companion check shows underprediction in a flat-heavy setting. The direction in other settlement patterns is unknown.
\item \textbf{External validity.} Available national comparisons use published totals or quantiles. Examples and parity checks use synthetic inputs. The direction of disagreement with observed contact is unknown.
\item \textbf{Computational scale.} The two existing presets differ in population size, employer limits and geography. They can create different edge and two-step path counts per person. Extrapolation to a nationwide run could underestimate cost in settings with larger dense groups. The amount is unknown.
\end{itemize}

\section{Conclusion}

A register-derived network consists of a population input and the choices used to make its ties. The \builtnet{} construction cards expose those choices and allow an analyst to rerun specific alternatives with the same memberships. The tests support the implementation on the checked inputs. The worked examples show how caps, geometry, age eligibility and overlap shape the resulting summaries. The unresolved Danish disagreement and the secondary-school prediction failure show where source documentation or better input is still needed.
